\DocumentMetadata{
  pdfversion=2.0,
  pdfstandard=ua-2,
  lang=en-US,
  tagging=on,
  tagging-setup={math/setup=mathml-SE,tabsorder=structure}
}
\documentclass[11pt,letterpaper]{article}
\usepackage[margin=0.68in,includefoot]{geometry}
\usepackage{newcomputermodern}
\usepackage{amsmath,amscd,mathtools}
\providecommand{\square}{\mdlgwhtsquare}
\usepackage[hidelinks]{hyperref}
\hypersetup{
  pdftitle={Analysis First-Year Exam, August 2024, Part 2},
  pdfauthor={University of Florida Department of Mathematics},
  pdfsubject={Graduate mathematics examination},
  pdfdisplaydoctitle=true,
  bookmarksopen=true
}
\setcounter{secnumdepth}{0}
\setlength{\parindent}{0pt}
\setlength{\parskip}{0.28em}
\setlength{\emergencystretch}{3em}
\setlength{\leftmargini}{2.15em}
\setlength{\leftmarginii}{2.65em}
\setlength{\leftmarginiii}{3em}
\renewcommand{\labelenumi}{\textbf{\theenumi.}}
\pagestyle{empty}
\raggedbottom
\allowdisplaybreaks
\newenvironment{examproblems}[1]{%
  \begin{enumerate}%
  \setcounter{enumi}{#1}%
  \setlength{\topsep}{0.35em}%
  \setlength{\partopsep}{0pt}%
  \setlength{\itemsep}{0.48em}%
  \setlength{\parsep}{0.18em}%
}{\end{enumerate}}
\newenvironment{examsubparts}{%
  \begin{enumerate}%
  \setlength{\topsep}{0.18em}%
  \setlength{\partopsep}{0pt}%
  \setlength{\itemsep}{0.16em}%
  \setlength{\parsep}{0.1em}%
}{\end{enumerate}}
\newenvironment{examinstructions}{%
  \par\begingroup\itshape
}{%
  \par\endgroup\vspace{0.18em}
}
\makeatletter
\renewcommand\section{\@startsection{section}{1}{0pt}%
  {-1.25ex plus -.25ex minus -.1ex}%
  {0.55ex plus .1ex}%
  {\normalfont\large\bfseries}}
\renewcommand{\@maketitle}{%
  \newpage\null\vskip -1.2em%
  {\footnotesize\bfseries
    \href{https://www.ufl.edu/}{University of Florida}, \href{https://math.ufl.edu/}{Department of Mathematics}\par}%
  \vskip 0.18em%
  \tagstructbegin{tag=Title}%
    {\LARGE\bfseries\href{https://gma.math.ufl.edu/exams/analysis-first-year/}{Analysis First-Year Exam}, August 2024, Part 2\par}%
  \tagstructend%
  \vskip 0.35em%
  \tagmcbegin{artifact}%
    \hrule height 1pt%
  \tagmcend%
  \vskip 0.55em%
}
\makeatother
\title{Analysis First-Year Exam, August 2024, Part 2}
\author{}
\date{}
\begin{document}
\maketitle
\thispagestyle{empty}
\begin{examinstructions}
Answer FOUR questions in detail. State carefully any results used without proof.
\end{examinstructions}
\begin{examproblems}{0}
\item
This problem has two parts.
\begin{examsubparts}
\item[\textnormal{(i)}]
The function \(f:[-1,1]\to\mathbb R\) is Riemann integrable and \(F:[-1,1]\to\mathbb R\) is defined by the rule 
\[
F(t)=\int_0^t f(u)\,du.
\]
 Prove that~\(F\) is continuous at~\(0\).
\item[\textnormal{(ii)}]
Does the same conclusion follow if~\(f\) is instead assumed to be Lebesgue integrable? Explain.
\end{examsubparts}
\item
Let \(A=\{a_n:n\in\mathbb N\}\) be a subset of the metric space~\(X\). For each \(n\in\mathbb N\) let \(f_n:X\to\mathbb R\) be continuous except perhaps at \(a_n\).
\begin{examsubparts}
\item[\textnormal{(i)}]
Prove that if \(f_n\to f\) uniformly on~\(X\) then the set of discontinuities of~\(f\) is finite or countably infinite.
\item[\textnormal{(ii)}]
Show by example that if \(f_n\to f\) pointwise on~\(X\) then~\(f\) can have uncountably many discontinuities.
\end{examsubparts}
\item
Let \(f:[0,1]\to\mathbb R\) be continuous; assume that 
\[
\int_0^1 f(t)t^n\,dt=0
\]
 for each integer \(n>1\). Prove that the function~\(f\) is constantly zero.
\item
Let \((\Omega,\mathcal A,\mu)\) be a measure space that is complete in the sense that each subset of every~\(\mu\)-null set lies in \(\mathcal A\). Let \(f:\Omega\to\mathbb R\) be measurable and let \(A\subseteq\Omega\) be~\(\mu\)-null. Prove that if \(g:\Omega\to\mathbb R\) is defined by changing the values of~\(f\) at points of~\(A\), then~\(g\) is measurable.
\item
On the measure space~\(\Omega\), let \((f_n:n\in\mathbb N)\) be a sequence of~\(\mu\)-integrable functions that converges uniformly to the function~\(f\). Prove that if \(\mu(\Omega)\) is finite then~\(f\) is~\(\mu\)-integrable, and show by example that~\(f\) can fail to be~\(\mu\)-integrable if \(\mu(\Omega)\) is infinite.
\end{examproblems}
\end{document}
